\section{Avaliação da resposta em linguagem natural}

Além das métricas sobre a consulta SQL, foi conduzida avaliação da \textbf{cadeia completa} do subsistema analítico: execução da consulta, geração da resposta ao cidadão pelo agente orquestrador e julgamento automático por outro modelo de linguagem (LLM-as-Judge).

\subsection{Protocolo}

Para cada item amostrado do experimento E1, o protocolo registrou: (i)~pergunta do usuário e SQL gerada; (ii)~resultado tabular executado no DuckDB; (iii)~resposta em linguagem natural produzida pelo orquestrador (GPT-4o); (iv)~avaliação do juiz (GPT-4.1-mini) em quatro dimensões Likert (adequação semântica, fidelidade aos dados, completude e clareza).

A métrica \textbf{JAR} (\textit{Judge Agreement Rate}) corresponde à proporção de respostas classificadas como \textit{corretas} ou \textit{parcialmente corretas} pelo juiz.

\subsection{Piloto ($n=5$)}

O piloto inicial confirmou a viabilidade do pipeline. Observou-se JAR de 100\% com EX de 0\%, reforçando que a métrica EX estrita subestima a utilidade percebida quando a resposta natural atende à pergunta com os dados retornados.

\subsection{Resultados estratificados ($n=30$)}

A amostra estratificada compreendeu 10 consultas fáceis, 10 médias e 10 difíceis. A Tabela~\ref{tab:jar} resume os resultados.

\begin{table}[htbp]
\centering
\caption{Avaliação LLM-as-Judge sobre respostas naturais (E1, $n=30$)}
\label{tab:jar}
\begin{tabular}{lrr}
\toprule
\textbf{Métrica} & \textbf{Valor} & \textbf{Observação} \\
\midrule
JAR & 100,00\% (30/30) & Respostas corretas ou parcialmente corretas \\
EX (mesma amostra) & 20,00\% (6/30) & Equivalência estrita SQL \\
Adequação semântica (média) & 4,87 / 5 & Escala Likert do juiz \\
Fidelidade aos dados (média) & 4,87 / 5 & \\
Completude (média) & 4,73 / 5 & \\
Clareza ao cidadão (média) & 4,77 / 5 & \\
Custo total da amostra & US\$ 0,30 & Orquestrador + juiz \\
\bottomrule
\end{tabular}
\end{table}

\begin{table}[htbp]
\centering
\caption{JAR e EX por dificuldade (amostra $n=30$)}
\label{tab:jar-diff}
\begin{tabular}{lrrr}
\toprule
\textbf{Dificuldade} & \textbf{$n$} & \textbf{JAR (\%)} & \textbf{EX (\%)} \\
\midrule
Fácil & 10 & 100,00 & 20,00 \\
Médio & 10 & 100,00 & 20,00 \\
Difícil & 10 & 100,00 & 20,00 \\
\bottomrule
\end{tabular}
\end{table}

\subsection{Interpretação}

O contraste entre EX (20\%) e JAR (100\%) na mesma amostra indica que, embora a formulação SQL frequentemente difira da referência manual, o orquestrador consegue comunicar resultados úteis e fidedignos aos dados executados. Essa evidência complementa as limitações da métrica EX e justifica a avaliação em duas camadas: sintaxe/resultado tabular e adequação da resposta ao usuário final.

Os registros completos de entrada e saída de cada etapa foram arquivados para auditoria e reprodutibilidade do experimento.
